% Draft sections for weight-optimization evaluation.
% Intended for insertion into Final_Report.tex under \section{Results}.
% Do NOT wrap in \documentclass / \begin{document}.

\subsection{Held-Out Validation of Weight Optimization}
\label{sec:holdout_weight}

The scoring weights (Equations~1--2) were originally optimized on the full feature matrix and evaluated on the same data, producing an AUC of 0.9685 recorded in \texttt{optimized\_weights.json}. Because Nelder-Mead explores the weight space against the very labels used for evaluation, that value is optimistically biased by construction. To determine whether the learned weights generalize, we re-ran the optimizer under a stratified calibration/evaluation split.

\subsubsection{Protocol}

A wrapper script (\texttt{analysis/optimize\_weights\_holdout.py}) loads the cached edge table via the same loader used by the feature audit, applies a stratified 50/50 split (seed~42), and partitions the 364,802 masked edges into a calibration half (182,402 edges, 153 red-team) and an evaluation half (182,400 edges, 152 red-team). The optimizer is constructed on the calibration half only, runs Nelder-Mead to convergence on the same five features used in the original pipeline (\texttt{is\_ntlm}, \texttt{source\_fan\_out}, \texttt{dst\_in\_degree}, \texttt{is\_network\_logon}, \texttt{dst\_fan\_out\_ratio}), and the resulting weights are applied to the evaluation half without further modification.

\subsubsection{Result: Optimization Generalizes}

Table~\ref{tab:holdout_optimizer} reports the calibration and evaluation AUCs.

\begin{table}[h]
\centering
\begin{tabular}{lc}
\hline
Quantity & Value \\
\hline
Calibration AUC & 0.968865 \\
Eval AUC (held-out) & \textbf{0.968070} \\
Calibration $-$ Eval gap & $+$0.000796 \\
Nelder-Mead iterations (calibration) & 143 \\
\hline
\end{tabular}
\caption{Held-out validation of the Nelder-Mead weight optimizer. Calibration and evaluation AUCs agree to within 0.0008, indicating no measurable overfitting at five free parameters over roughly 365{,}000 edges.}
\label{tab:holdout_optimizer}
\end{table}

Calibration and evaluation AUCs agree to within 0.0008, well inside the variation expected from 152/153 red-team edges per split. With five free parameters against 305 red-team edges (approximately 0.016 parameters per positive sample), the hypothesis class is extremely under-parameterized, leaving essentially no capacity to overfit. The 0.968 figure in the original optimization trace is not an in-sample artifact; the same weights would land near 0.968 on an unseen sample drawn from the same distribution. The AUC improvement from manually-tuned weights ($\approx$0.954) to optimized weights (0.968) is therefore defensible with no held-out penalty.


\subsection{Logistic-Regression Baseline Comparison}
\label{sec:lr_baseline}

To test whether the bespoke Nelder-Mead optimizer adds value over a standard supervised linear method, we fit logistic regression (L2 penalty, balanced class weights) on the calibration half and evaluated on the held-out half, with identical feature preprocessing: the same percentile-rank transform applied to \texttt{edge\_rarity} and \texttt{protocol\_rarity}, followed by standardization on calibration-half statistics.

\begin{table}[h]
\centering
\begin{tabular}{lcccc}
\hline
Method & Cal AUC & Eval AUC & Gap & Wall time \\
\hline
Nelder-Mead optimizer   & 0.9689 & 0.9681 & $+$0.0008 & $\approx$25\,s \\
Logistic regression (L2) & 0.9738 & \textbf{0.9733} & $+$0.0005 & $\approx$0.2\,s \\
\hline
\end{tabular}
\caption{Nelder-Mead optimizer versus logistic regression on the same held-out split. The 0.005 eval-AUC gap in favor of LR is within single-seed noise.}
\label{tab:lr_comparison}
\end{table}

Logistic regression and the Nelder-Mead optimizer are essentially equivalent on this data (Table~\ref{tab:lr_comparison}). The small numerical difference favoring LR is likely attributable to its intercept term and L2 regularization. The specific choice of supervised linear method does not materially affect the result: Nelder-Mead, logistic regression, and most other supervised linear methods on these five features all land near AUC~0.97. The optimizer's value lies in producing interpretable weights aligned with the scoring formula, not in outperforming off-the-shelf alternatives.


\subsection{Tabular versus Graph-Derived Feature Ablation}
\label{sec:tabular_graph_ablation}

The Nelder-Mead-versus-LR comparison varies only the learner; both methods use the same five features. To isolate whether the features themselves, and specifically the graph-derived ones, carry the discriminative signal, we ran a feature ablation that holds the learner constant (logistic regression, same held-out split, seed~42) and varies the feature set.

\subsubsection{Feature Grouping}

Features were split into two groups by whether they require graph topology to compute. \emph{Pure tabular} features are computable from per-edge attributes alone: authentication flags, edge rarity (a function of edge weight), protocol rarity, byte-per-packet ratio, duration z-score, and the unusual-destination-port flag. \emph{Graph-derived} features require traversing the graph: node degrees in both directions, fan-out ratios, node-level temporal features (inter-arrival statistics, burst score, active duration), and destination-node betweenness centrality.

\subsubsection{Result: Graph-Derived Features Dominate}

\begin{table}[h]
\centering
\begin{tabular}{lccc}
\hline
Feature set & \# features & Cal AUC & Eval AUC \\
\hline
A. Pure tabular only    & 9  & 0.9609 & 0.9562 \\
B. Graph-derived only   & 17 & 0.9893 & \textbf{0.9891} \\
C. Combined (A $+$ B)  & 26 & 0.9981 & 0.9922 \\
\hline
\end{tabular}
\caption{Feature ablation under logistic regression. Adding graph-derived features on top of pure-tabular raises eval AUC by $+$0.0359; adding pure-tabular on top of graph-derived raises eval AUC by only $+$0.0030.}
\label{tab:tabular_graph}
\end{table}

Table~\ref{tab:tabular_graph} shows that graph-derived features carry roughly an order of magnitude more discriminative signal than pure-tabular features on this dataset. Adding graph-derived features on top of the pure-tabular set raises eval AUC by $+$0.0359, while adding pure-tabular features on top of the graph-derived set raises eval AUC by only $+$0.0030.

This result directly supports the graph-based methodology. The five features used in the original scoring function included three graph-derived columns (\texttt{source\_fan\_out}, \texttt{dst\_in\_degree}, \texttt{dst\_fan\_out\_ratio}); the LR-versus-optimizer comparison therefore tested learner-versus-learner on graph features, not graph-versus-no-graph. The proper ablation above isolates the feature contribution and confirms that graph-derived features do the bulk of the discriminative work.


\subsection{Graph Feature Extensions}
\label{sec:graph_extensions}

The ablation in Section~\ref{sec:tabular_graph_ablation} shows that local graph features (degrees, fan-out, node-level temporal patterns) carry most of the signal. A natural follow-up is whether multi-hop graph reasoning, features that require traversing more than one hop, adds further improvement. We swept five candidate graph feature groups, added one at a time on top of the five-feature LR baseline.

\begin{table}[h]
\centering
\begin{tabular}{lccc}
\hline
Feature group & Cal AUC & Eval AUC & $\Delta$ vs base \\
\hline
(base: 5 features only)                         & 0.9738 & 0.9733 & --- \\
Standard PageRank                                & 0.9739 & 0.9734 & $+$0.0002 \\
Personalized PageRank (seed: \texttt{C17693})   & \textbf{0.9964} & \textbf{0.9956} & \textbf{$+$0.0224} \\
k-core decomposition                            & 0.9773 & 0.9771 & $+$0.0039 \\
Louvain community                                & 0.9779 & 0.9779 & $+$0.0046 \\
Jaccard $+$ Adamic-Adar similarity              & 0.9737 & 0.9733 & $+$0.0001 \\
All five groups stacked                          & 0.9968 & 0.9965 & $+$0.0232 \\
\hline
\end{tabular}
\caption{Quick-win graph-feature sweep on top of the five-feature LR baseline. Louvain community features include a cross-community flag and community-size features.}
\label{tab:graph_sweep}
\end{table}

\subsubsection{Personalized PageRank}

Personalized PageRank seeded at the known compromised host (\texttt{C17693}) adds far more than any other feature in the sweep (Table~\ref{tab:graph_sweep}). This is a real signal, but it carries a methodological caveat: the feature uses the identity of \texttt{C17693} as an input, and in a production deployment that identity would have to be discovered first. The result supports a \emph{conditional} claim: given a seed of confirmed-compromised hosts, structural propagation lifts eval AUC from 0.973 to 0.996. Carefully framed as a post-detection scoping mechanism, this is publishable; presented without qualification, it would overstate what the method can achieve under cold-start conditions.

\subsubsection{Label-Free Multi-Hop Features}

The label-free graph features, k-core decomposition and Louvain community membership, each add a modest but consistent improvement of roughly 0.004 eval AUC over the local-graph baseline. These features require no prior information about attackers and provide direct support for the claim that multi-hop graph structure carries detectable signal beyond local features. Standard PageRank and edge-similarity scores (Jaccard, Adamic-Adar) contribute essentially nothing on this dataset.

Stacking all five feature groups reaches eval AUC 0.9965, but the gain is dominated by personalized PageRank. With personalized PageRank removed from the stack, the remaining label-free graph features push eval AUC to approximately 0.978.


\subsection{Decomposing the Detection Improvement}
\label{sec:weight_synthesis}

The four results above isolate different axes of the detection pipeline. Table~\ref{tab:synthesis_axes} summarizes what each comparison tests.

\begin{table}[h]
\centering
\begin{tabular}{lp{2.2cm}p{2.0cm}p{2.2cm}}
\hline
Comparison & Held constant & What varies & Measures \\
\hline
Result~1 & Optimizer, features, split & Which half is evaluated & Overfitting? \emph{No.} \\
Result~2 & Features, split, supervision & Linear method & Optimizer choice? \emph{No measurable effect.} \\
Result~3 & Learner, supervision, split & Feature set & Graph features? \emph{Yes, $+$0.036.} \\
Result~4 & Learner, base features, split & Graph feature group & Multi-hop features? \emph{Modestly, label-free.} \\
\hline
\end{tabular}
\caption{Summary of what each held-out comparison isolates.}
\label{tab:synthesis_axes}
\end{table}

Combining these results with the existing one-class baseline comparisons in Table~\ref{tab:methods_lanl} decomposes the improvement from na\"ive detection to the best achievable AUC into additive contributions (Table~\ref{tab:auc_budget}). Isolation Forest on graph-derived features (AUC~0.94) versus logistic regression on the same features (AUC~0.99) shows a roughly 0.05~AUC improvement attributable to having labels at training time. Graph-derived local features contribute a further $+$0.036 over pure-tabular features under the same supervised learner. Multi-hop label-free features add $+$0.004 to $+$0.007 on top. The largest single contribution is the graph-derived features themselves, not the choice of optimizer, not the multi-hop reasoning, and not the path-boost mechanism.

\begin{table}[h]
\centering
\begin{tabular}{lc}
\hline
Contribution & Approx.\ AUC $\Delta$ \\
\hline
Pure-tabular features under one-class learning & (lower bound) \\
Pure-tabular features under supervised LR      & $+$0.05 (inferred) \\
Graph-derived local features under supervised LR & $+$0.036 \\
Multi-hop graph features (k-core, communities)  & $+$0.004--0.007 \\
Known-attacker propagation (personalized PageRank) & $+$0.022 (conditional) \\
\hline
\end{tabular}
\caption{Approximate AUC budget from the worst defensible baseline to the most aggressive supervised graph-aware system. The largest single addition is the graph-derived features themselves.}
\label{tab:auc_budget}
\end{table}

Three claims emerge from this evaluation with clear empirical backing. First, graph-derived features carry the bulk of the discriminative signal for lateral-movement detection on LANL-2015: substituting pure-tabular features for graph-derived features under the same supervised linear learner drops eval AUC by roughly 0.036. Second, supervised learning is a separate, additive contribution: logistic regression on graph-derived features (eval AUC 0.99) versus one-class baselines on the same features (Isolation Forest at AUC 0.94) yields a roughly 0.05~AUC improvement attributable to labels at training time. Third, among supervised linear methods on the same features, the specific algorithm matters little: Nelder-Mead direct AUC maximization and logistic regression agree to within single-seed noise. The optimizer's value is in producing interpretable weights aligned with the scoring formula, not in outperforming standard alternatives.
